\section{Introduction}\label{sec:intro}

Cyber risk has moved from the information-technology department to the agenda of financial-stability authorities. The International Monetary Fund devotes a chapter of its \emph{Global Financial Stability Report} to the macrofinancial consequences of cyber incidents \citep{imf2024}; the European Union's Digital Operational Resilience Act brings critical information and communication technology providers under supervisory oversight \citep{dora2022}; and a growing academic literature asks whether a cyber attack on a large bank could propagate like a run or a payment-system failure \citep{duffie2019,kashyap2019,eisenbach2022}. Behind these developments lies a deceptively simple question that supervisors, insurers and boards keep asking: \emph{when} will a catastrophic cyber event hit the banking system, and \emph{how large} will it be?

Two empirical regularities frame any answer. The first concerns \emph{where} attackers strike. The largest thefts of the past decade have occurred in crypto-asset venues rather than in the core of the banking system. Cross-chain bridges and custodial exchanges have lost hundreds of millions to more than a billion dollars in single incidents, including the theft of roughly \$1.5 billion from a major exchange in February 2025 \citep{fbi2025,chainalysis2025}. The best-known attack on the banking core, the 2016 compromise of a central bank's payment-messaging environment, issued fraudulent instructions worth close to \$1 billion, yet most were blocked and the amount ultimately lost was an order of magnitude smaller. Banks hold vastly larger balances than any crypto venue, so this pattern calls for an explanation. The second regularity concerns \emph{how large} losses can be. Cyber losses are heavy-tailed \citep{maillart2010,edwards2016,eling2019}, so the quantities of policy interest, such as the 1-in-100-year or 1-in-1,000-year loss, lie far beyond the range of any observed record. They can only be reached by extrapolation, and extrapolation from a few decades of data is precisely where statistical confidence is most likely to be misplaced.

This paper studies both regularities, and the channel that links them to systemic risk, within one reproducible framework. We ask three questions. \emph{Why} do rational attackers prefer venues that hold less money? \emph{How confidently} can one state the return period of a catastrophic banking-sector loss from realistic data? And \emph{how} does a large cyber loss at one institution propagate through interbank exposures? \Cref{fig:design} summarizes the design: an attacker model determines which venues are attractive; a loss process with a known, planted tail generates incident and annual losses; an inference layer estimates extreme return levels from short records with three interval methods; an audit layer measures how often those intervals contain the truth; and a contagion layer propagates large losses through an interbank network.

\begin{figure}[!tbp]
\centering
\begin{tikzpicture}[
  box/.style={draw=inkX!70, rounded corners=3pt, fill=#1, align=center, font=\footnotesize, inner sep=5pt, minimum height=1.55cm, text width=3.25cm},
  box/.default={white},
  arr/.style={-{Stealth[length=2.2mm]}, draw=inkX!70, line width=0.6pt}]
\node[box] (att) {\textbf{Attacker incentives}\\[2pt] venue choice and the finality lever (\S\ref{sec:attacker})};
\node[box, right=0.6cm of att] (loss) {\textbf{Loss process}\\[2pt] Poisson frequency and a planted tail, $\xi = \xiTrue$ (\S\ref{sec:lda})};
\node[box=violetX!8, right=0.6cm of loss] (inf) {\textbf{Inference}\\[2pt] POT--GPD fit to an $R$-year record; three intervals (\S\ref{sec:evt})};
\node[box=violetX!8, right=0.6cm of inf] (aud) {\textbf{Audit}\\[2pt] coverage against the known truth (\S\ref{sec:design})};
\node[box=hotX!7, anchor=north, text width=8.4cm, minimum height=1.0cm] (net) at ($(inf.south)!0.5!(aud.south) + (0,-0.55)$) {\textbf{Contagion}: linear DebtRank on an interbank network, with a phase transition at $\rho(\Lambda)=1$ (\S\ref{sec:contagion})};
\draw[arr] (att) -- (loss);
\draw[arr] (loss) -- (inf);
\draw[arr] (inf) -- (aud);
\draw[arr] (loss.south) |- (net.west);
\end{tikzpicture}
\caption{Research design. The loss process carries a planted tail index, so every inference made from a short record can be checked against a truth computed from \TruthYears{} simulated years. Shaded boxes contain the paper's main methodological contribution.}
\label{fig:design}
\end{figure}

Our central finding is that standard practice is overconfident in a specific and dangerous way. In a controlled Monte Carlo experiment in which the truth is known, nominal 95\% intervals for the 1-in-1,000-year loss built with the delta method or the parametric bootstrap contain the truth only \covBThirtyDeltaK\% and \covBThirtyBootstrapK\% of the time when estimated from 30-year records generated by a realistic loss process. The failures are entirely one-sided: \belowBThirtyDeltaK\% of delta intervals lie wholly \emph{below} the true loss and none lie above it. The reason is that the point estimate itself is severely biased downward in short records: the median estimate is only \medRatioBThirtyK{} times the true 1-in-1,000-year loss, because a short record rarely contains the rare events that reveal how heavy the tail is. Intervals that are symmetric on the log scale around such an estimate cannot reach the truth. Profile-likelihood intervals, which are free to be strongly right-skewed, raise coverage to \covBThirtyProfileK\% in the same setting and to between \covAHundredProfileK\% and \covAThreeHundredProfileK\% once the record contains 50 or more exceedances. They also roughly halve the proper interval score of \citet{gneiting2007}, so their additional width is not mere caution but better calibrated uncertainty. The price is width, exactly where data are thin: with 15 exceedances the median upper limit of a profile interval is \ratioAThirtyProfileK{} times its lower limit. Such a ratio is not a defect of the method; it is an honest statement of how little a few decades of data can say about a millennial loss.

The paper makes five contributions. \emph{First}, we formalize attacker venue selection as a rational-choice problem in the tradition of \citet{becker1968} and show that, because retained value is a product of four factors, the log payoff gap between two venues decomposes \emph{exactly} into factor contributions. With illustrative parameters, finality explains \shareFinality\% of the gap between a bank core and a cross-chain bridge, and a bank's recovery capability is the single most effective deterrent. \emph{Second}, we build a loss-distribution model whose tail index is known by construction, which turns extreme-value inference into an auditable experiment, and we derive a second-order compound-sum correction, $\psi^{S}_{T} \approx \psi_{T} + \E[S]$, that converts incident-level return levels into the annual-aggregate levels that matter for solvency. \emph{Third}, we run a coverage study with \Reps{} independent records in each of twelve design cells and compare delta-method, parametric-bootstrap and profile-likelihood intervals using Wilson confidence bands, paired McNemar tests and proper interval scores. \emph{Fourth}, we show that what profile likelihood cannot repair is threshold bias: with long records, a threshold placed near a structural kink in the severity distribution inflates the fitted shape and turns the remaining errors into overestimates. \emph{Fifth}, we model interbank contagion with linear DebtRank \citep{bardoscia2015}, verify its closed form numerically, and show a phase transition at a critical exposure of \critExp{} times equity, beyond which a cyber hit on a shared technology vendor destroys \scVendHighShare\% of system equity. Systemic importance in our network tracks interconnection rather than size: a periphery bank with \pOneAssets{} of assets outranks core banks \pOneRatioLo{} to \pOneRatioHi{} times larger.

All parameters are illustrative, chosen to be plausible rather than estimated. The contribution is therefore methodological: the behaviour of interval estimators, the size of the compound-sum correction and the existence of the phase transition depend on the tail index, the number of exceedances and the network spectrum, not on the dollar values. The conclusions transfer to any heavy-tailed loss process with a comparable shape and comparable data scarcity; the dollar figures do not. We return to calibration in \cref{sec:limits}.

The rest of the paper is organized as follows. \Cref{sec:literature} reviews related work. \Cref{sec:framework} presents the attacker model, the loss process, the inference methods and the contagion model. \Cref{sec:design} describes the experimental design. \Cref{sec:results} reports results, \cref{sec:discussion} discusses implications for supervision and market design, \cref{sec:limits} states limitations, and \cref{sec:conclusion} concludes. The appendix contains derivations, the full coverage table and implementation details.
